% NB15b -- Caldera de domo con circulacion natural (modelo de Astrom-Bell 2000): hogar (gas de escape de la TG
% o biomasa EFB/PMF), elevadores, bajantes, domo con vapor bajo la superficie; variables de estado y el
% fenomeno de "shrink & swell". Reemplaza "SERPENTIN.png" de la entrega PlantaTermoelectrica.
\documentclass[tikz,border=4mm]{standalone}
\usepackage{amsmath,amssymb}
\usetikzlibrary{arrows.meta,calc,patterns,decorations.pathmorphing}
\begin{document}
% Estilo comun de los diagramas TikZ del modulo VCR (NB10b) -- paleta del proyecto
\definecolor{manivela}{HTML}{2A78D6}
\definecolor{inferior}{HTML}{4A3AA7}
\definecolor{control}{HTML}{1BAF7A}
\definecolor{superior}{HTML}{EB6834}
\definecolor{tinta}{HTML}{3A3A38}
\definecolor{metal}{HTML}{C3C2B7}
\definecolor{perdida}{HTML}{E34948}
\definecolor{fuerza}{HTML}{C0392B}
\tikzset{
  barra/.style={line width=3.2pt, line cap=round},
  pin/.style={circle, fill=white, draw=tinta, line width=0.9pt, inner sep=1.6pt},
  fijo/.style={circle, fill=tinta, inner sep=2.6pt},
  vec/.style={-{Stealth[length=2.6mm]}, line width=0.9pt},
  F/.style={-{Stealth[length=3mm]}, line width=1.2pt, fuerza},
  cota/.style={{Stealth[length=1.8mm]}-{Stealth[length=1.8mm]}, line width=0.5pt, tinta},
  lbl/.style={font=\small, text=tinta},
}
\newcommand{\tierra}[1]{\fill[pattern=north east lines, pattern color=tinta] ($(#1)+(-5,-3)$) rectangle ($(#1)+(5,-6)$);
  \draw[tinta, line width=0.8pt] ($(#1)+(-5,-3)$) -- ($(#1)+(5,-3)$);}

\definecolor{frio}{HTML}{2A78D6}
\definecolor{vapor}{HTML}{4A3AA7}
\definecolor{caliente}{HTML}{E34948}
\definecolor{tibio}{HTML}{EB6834}
\begin{tikzpicture}[x=1cm,y=1cm,
  flujo/.style={-{Stealth[length=2.6mm]}, line width=1.6pt}]
  \node[lbl, font=\small\bfseries, anchor=west] at (-1.2,9.0) {(a) Circuito de circulaci\'on natural y estados del modelo};
  % hogar
  \draw[tinta, line width=1.2pt, fill=tibio!10] (0,0) rectangle (5,5.6);
  \fill[pattern=crosshatch, pattern color=tibio] (0.3,0.15) rectangle (4.7,0.45);
  \node[lbl, font=\scriptsize] at (2.5,0.75) {parrilla (EFB/PMF) o gas de la TG};
  \foreach \x in {1.2,2.5,3.8} \draw[tibio, line width=1.2pt, decorate, decoration={snake, amplitude=1.5pt, segment length=7pt}] (\x,1.0) -- (\x,2.6);
  \node[lbl, font=\footnotesize, tibio!80!black, align=center] at (2.5,3.2) {$\dot Q$ (radiaci\'on\\+ convecci\'on)};
  % elevadores (paredes de agua)
  \foreach \x in {5.15,5.35,5.55} \draw[vapor!70, line width=1.6pt] (\x,0.4) -- (\x,6.8);
  \foreach \y in {1.5,2.6,3.7,4.8,5.9} \fill[white] (5.35,\y) circle (0.07) (5.15,\y+0.4) circle (0.06) (5.55,\y+0.25) circle (0.05);
  \node[lbl, font=\footnotesize, vapor, rotate=90] at (5.95,3.2) {elevadores $V_r$: mezcla, calidad $\alpha_r$};
  % bajante
  \draw[frio, line width=3pt] (9.6,6.8) -- (9.6,0.2) -- (5.35,0.2) -- (5.35,0.4);
  \node[lbl, font=\footnotesize, frio, rotate=90] at (9.95,3.2) {bajantes $V_{dc}$: l\'iquido saturado};
  \draw[-{Stealth[length=2.4mm]}, frio, line width=1pt] (9.25,4.0) -- (9.25,3.0);
  \draw[-{Stealth[length=2.4mm]}, vapor, line width=1pt] (6.3,2.0) -- (6.3,3.0);
  \node[lbl, font=\scriptsize, align=center] at (7.6,-0.75) {circulaci\'on: $q_{dc}=\sqrt{\dfrac{2\rho_w A_{dc}^2(\rho_w-\rho_s)g\,\bar\alpha_v V_r}{k}}$};
  % domo
  \draw[tinta, line width=1.3pt, fill=white] (5.0,6.5) rectangle (10.2,8.4);
  \fill[frio!30] (5.05,6.55) rectangle (10.15,7.5);
  \draw[frio!80!black, line width=1.2pt] (5.05,7.5) -- (10.15,7.5);
  \foreach \p in {(5.6,7.0),(6.1,6.8),(6.7,7.15),(7.2,6.85),(5.9,7.3),(6.5,6.7)} \draw[vapor, line width=0.6pt, fill=white] \p circle (0.08);
  \node[lbl, font=\scriptsize, vapor, align=center] at (7.0,8.0) {vapor bajo la superficie $V_{sd}$};
  \draw[cota] (10.5,6.5) -- (10.5,7.5); \node[lbl, font=\footnotesize, right] at (10.55,7.0) {nivel $\ell$};
  \draw[flujo, vapor] (9.0,8.4) -- (9.0,9.1) -- (11.2,9.1) node[right, font=\footnotesize] {$q_s$ al sobrecalentador};
  \draw[flujo, frio] (12.2,7.9) -- (10.2,7.9) node[pos=0.35, above, font=\footnotesize] {$q_f,\;h_f$};
  \node[lbl, font=\footnotesize, anchor=west] at (5.1,8.75) {domo $V_d$, \'area $A_d$};
  % metal
  \node[lbl, font=\scriptsize, align=left, anchor=north west] at (-1.2,-1.5) {metal: $m_t$ (total), $m_r$ (elevadores), $c_p=650$ J/kg K, $T_m=T_{sat}(p)$};
  % ================= estados y ecuaciones =================
  \node[draw=tinta, rounded corners=2pt, fill=white, font=\footnotesize, align=left, anchor=north west] at (13.6,8.9)
    {\textbf{Estados}: $p$ (domo), $V_{wt}$ (agua total), $\alpha_r$ (calidad a la salida\\de los elevadores), $V_{sd}$ (vapor bajo la superficie)\\[4pt]
     $e_{11}\dot V_{wt}+e_{12}\dot p=q_f-q_s$ \hfill (masa global)\\[2pt]
     $e_{21}\dot V_{wt}+e_{22}\dot p=Q+q_fh_f-q_sh_s$ \hfill (energ\'ia global)\\[2pt]
     $e_{32}\dot p+e_{33}\dot\alpha_r=Q-\alpha_rh_c\,q_{dc}$ \hfill (elevadores)\\[2pt]
     $e_{42}\dot p+e_{43}\dot\alpha_r+e_{44}\dot V_{sd}=\dfrac{\rho_s}{T_d}(V_{sd}^0-V_{sd})+\dfrac{h_f-h_w}{h_c}q_f$\\[4pt]
     $\ell=\dfrac{V_{wd}+V_{sd}}{A_d}$,\qquad $e_{ij}=f(p,\alpha_r;\;\rho,h,\partial_p\rho,\partial_p h)$ CoolProp};
  % ================= shrink & swell =================
  \begin{scope}[shift={(13.6,0)}]
    \node[lbl, font=\small\bfseries, anchor=west] at (0,4.3) {(b) Respuesta inversa (\emph{swell}) a un escal\'on de $q_s$};
    \draw[-{Stealth[length=2mm]}, tinta] (0,0) -- (6.6,0) node[below, font=\scriptsize] {$t$};
    \draw[-{Stealth[length=2mm]}, tinta] (0,-0.2) -- (0,3.6);
    \draw[tinta!50, densely dotted] (0,1.6) -- (6.4,1.6);
    \draw[tinta, line width=0.8pt] (0,0.4) -- (1,0.4) -- (1,0.9) -- (6.4,0.9) node[right, font=\scriptsize] {$q_s$};
    \draw[caliente, line width=1.4pt] plot[smooth] coordinates {(0,1.6) (1,1.6) (1.3,2.5) (1.8,2.9) (2.6,2.6) (3.6,1.9) (4.8,1.1) (6.4,0.25)};
    \node[caliente, font=\scriptsize, anchor=west] at (1.9,3.15) {$\ell$ sube primero};
    \draw[frio, line width=1.2pt, dashed] plot[smooth] coordinates {(0,1.6) (1,1.6) (2.5,1.3) (4.0,0.95) (6.4,0.45)};
    \node[frio, font=\scriptsize, anchor=west] at (4.1,1.35) {masa de agua $\downarrow$};
    \node[lbl, font=\scriptsize, align=left, anchor=north west] at (0,-0.35)
      {$q_s\uparrow\Rightarrow p\downarrow\Rightarrow$ \emph{flash} en elevadores y domo $\Rightarrow V_{sd}\uparrow$\\
       $\Rightarrow\ell\uparrow$ aunque el inventario de agua baja: control de nivel de\\
       un elemento act\'ua al rev\'es $\Rightarrow$ control de 3 elementos ($\ell$, $q_s$, $q_f$)};
  \end{scope}
\end{tikzpicture}
\end{document}
